\documentclass[11pt]{article}
\usepackage[margin=1in]{geometry}
\usepackage{amsmath,amssymb,amsthm,mathtools,bm}
\usepackage{xcolor}
\usepackage{listings}
\usepackage[colorlinks=true,linkcolor=blue!55!black,urlcolor=blue!55!black]{hyperref}
\usepackage{enumitem}

\theoremstyle{plain}
\newtheorem{lemma}{Lemma}
\newtheorem{theorem}{Theorem}
\newtheorem{proposition}{Proposition}
\theoremstyle{definition}
\newtheorem{remark}{Remark}

\newcommand{\R}{\mathbb{R}}
\newcommand{\E}{\mathbb{E}}
\newcommand{\Prob}{\mathbb{P}}
\newcommand{\N}{\mathcal{N}}
\newcommand{\D}{\mathcal{D}}
\newcommand{\C}{\mathcal{C}}
\newcommand{\K}{\mathbf{K}}
\newcommand{\I}{\mathbf{I}}
\newcommand{\x}{\mathbf{x}}
\newcommand{\y}{\mathbf{y}}
\newcommand{\bphi}{\boldsymbol{\phi}}
\newcommand{\btheta}{\boldsymbol{\theta}}
\newcommand{\bSigma}{\boldsymbol{\Sigma}}
\newcommand{\hy}{\hat{y}}
\newcommand{\indic}{\mathbb{I}}
\newcommand{\sn}{\sigma_n}
\newcommand{\code}[1]{\texttt{#1}}

\definecolor{codebg}{RGB}{247,247,249}
\definecolor{codekw}{RGB}{0,90,160}
\definecolor{codecm}{RGB}{110,120,130}
\definecolor{codestr}{RGB}{160,60,60}
\lstset{
  language=Python, basicstyle=\ttfamily\footnotesize,
  backgroundcolor=\color{codebg}, keywordstyle=\color{codekw},
  commentstyle=\color{codecm}\itshape, stringstyle=\color{codestr},
  showstringspaces=false, breaklines=true, breakatwhitespace=false,
  columns=fullflexible, frame=single, framerule=0.2pt, rulecolor=\color{gray!40},
  xleftmargin=6pt, xrightmargin=4pt, aboveskip=6pt, belowskip=4pt,
  literate={~}{{$\sim$}}1
}

\title{\vspace{-1.2cm}\bfseries Online Scalable Gaussian Processes with\\Conformal Prediction: A Rigorous Derivation Companion\\[2pt]\large with the implementing code}
\author{Derivation notes on Xu, Lu \& Giannakis (ICASSP 2025; arXiv:2410.05444)}
\date{}

\begin{document}
\maketitle
\vspace{-0.7cm}

\noindent\textbf{Scope and conventions.} This document follows the paper section by section and fills in, with rigour, the derivations the paper compresses. Displayed equations are transcribed as in the manuscript; the surrounding text is a paraphrase, with only short quoted phrases. Each result is paired with the function in \code{Ensemble Guassian Process.py} that implements it. In the code, \code{sigma\_f} denotes the kernel magnitude $\sigma_\theta$, \code{sigma\_n} the observation-noise standard deviation $\sn$, and the array \code{V} stores the $D$ sampled spectral frequencies, one per column. Throughout, $\{(\x_t,y_t)\}_{t\ge1}$ is a stream in $\R^d\times\R$, and $\D_t=\{\x_\tau,y_\tau\}_{\tau\le t}$.

\tableofcontents
\bigskip

%======================================================================
\section{Setup and the observation model}
The learner sees inputs $\x_t$ and scalar outputs $y_t$ generated by
\begin{equation}
y_\tau = f(\x_\tau) + n_\tau,\qquad n_\tau\sim\N(0,\sn^2)\ \text{i.i.d.},
\label{eq:obs}
\end{equation}
with $f$ unknown. Two goals are in tension: (i) a \emph{point} predictor $\hy_t(\x)$, and (ii) a \emph{set} predictor $\C_t(\x)$ that contains the true $y$ with a prescribed probability $1-\alpha$. The Gaussian process supplies both under its modelling assumptions; conformal prediction repairs the set's coverage when those assumptions fail. Sections~\ref{sec:gp}--\ref{sec:cp} recall the two ingredients; Sections~\ref{sec:rff}--\ref{sec:theorem} build the online scalable predictor and its guarantee.

%======================================================================
\section{Gaussian-process regression (\S II.A)}\label{sec:gp}
Place a zero-mean GP prior $f\sim\mathcal{GP}(0,\kappa(\x,\x'))$ with positive-definite kernel $\kappa$. For inputs $\x_1,\dots,\x_t$ the vector $\mathbf f_t=[f(\x_1),\dots,f(\x_t)]^\top$ is Gaussian, $p(\mathbf f_t\mid\x_{1:t})=\N(\mathbf f_t;\mathbf 0,\K_t)$ with $[\K_t]_{ij}=\kappa(\x_i,\x_j)$. Under \eqref{eq:obs} the labels satisfy $\y_t\mid\mathbf f_t\sim\N(\mathbf f_t,\sn^2\I_t)$, so marginally $\y_t\sim\N(\mathbf 0,\K_t+\sn^2\I_t)$. Writing $\mathbf k_t(\x)=[\kappa(\x_1,\x),\dots,\kappa(\x_t,\x)]^\top$, the joint law of the observed labels and a new (noisy) test label is Gaussian (paper eq.~1):
\begin{equation}
\begin{bmatrix}\y_t\\ y\end{bmatrix}\sim\N\!\left(\mathbf 0,\ \begin{bmatrix}\K_t+\sn^2\I_t & \mathbf k_t(\x)\\ \mathbf k_t^\top(\x) & \kappa(\x,\x)+\sn^2\end{bmatrix}\right).
\label{eq:joint}
\end{equation}

\begin{lemma}[Gaussian conditioning]\label{lem:cond}
If $\begin{bsmallmatrix}\mathbf a\\ \mathbf b\end{bsmallmatrix}\sim\N\!\big(\begin{bsmallmatrix}\boldsymbol\mu_a\\ \boldsymbol\mu_b\end{bsmallmatrix},\begin{bsmallmatrix}\mathbf A&\mathbf C\\ \mathbf C^\top&\mathbf B\end{bsmallmatrix}\big)$ with $\mathbf A\succ0$, then
$\mathbf b\mid\mathbf a\sim\N\big(\boldsymbol\mu_b+\mathbf C^\top\mathbf A^{-1}(\mathbf a-\boldsymbol\mu_a),\ \mathbf B-\mathbf C^\top\mathbf A^{-1}\mathbf C\big).$
\end{lemma}
\begin{proof}
Form $\mathbf z=\mathbf b-\mathbf C^\top\mathbf A^{-1}\mathbf a$. Then $\mathrm{Cov}(\mathbf z,\mathbf a)=\mathbf C^\top-\mathbf C^\top\mathbf A^{-1}\mathbf A=\mathbf 0$, so $\mathbf z\perp\mathbf a$ (jointly Gaussian $\Rightarrow$ independence). Hence $\mathbf b\mid\mathbf a$ has mean $\boldsymbol\mu_b+\mathbf C^\top\mathbf A^{-1}(\mathbf a-\boldsymbol\mu_a)$ and covariance $\mathrm{Cov}(\mathbf z)=\mathbf B-\mathbf C^\top\mathbf A^{-1}\mathbf C$.
\end{proof}

Applying Lemma~\ref{lem:cond} to \eqref{eq:joint} with $\mathbf a=\y_t$, $\mathbf b=y$, $\mathbf A=\K_t+\sn^2\I_t$, $\mathbf C=\mathbf k_t(\x)$, $\mathbf B=\kappa(\x,\x)+\sn^2$ gives the predictive pdf (paper eqs.~2--3):
\begin{align}
p(y\mid\D_t,\x)&=\N\big(y;\ \hy_t(\x),\ \sigma_t^2(\x)\big),\label{eq:pred}\\
\hy_t(\x)&=\mathbf k_t^\top(\x)\,(\K_t+\sn^2\I_t)^{-1}\y_t,\label{eq:pred-mean}\\
\sigma_t^2(\x)&=\kappa(\x,\x)-\mathbf k_t^\top(\x)(\K_t+\sn^2\I_t)^{-1}\mathbf k_t(\x)+\sn^2.\label{eq:pred-var}
\end{align}
The $\beta$-credible set is $K_t^\beta(\x)=[\hy_t(\x)-c_\beta\sigma_t(\x),\ \hy_t(\x)+c_\beta\sigma_t(\x)]$ with $c_\beta=\Phi^{-1}(\tfrac{1+\beta}{2})$; for $\beta=95\%$, $c_\beta\approx2$.

\begin{remark}[Cost and its consequence]
Equations \eqref{eq:pred-mean}--\eqref{eq:pred-var} require solving with $\K_t+\sn^2\I_t\in\R^{t\times t}$, an $O(t^3)$ factorisation that grows without bound in the stream. \textbf{No function forms $\K_t$}; Section~\ref{sec:rff} replaces the kernel by a finite feature map so that the same predictive becomes an $O((2D)^2)$ recursive update.
\end{remark}

%======================================================================
\section{Conformal prediction and the coverage guarantee (\S II.B)}\label{sec:cp}
Let $s_t(\x,y)$ be a \emph{non-conformity score}: larger means $y$ conforms worse to the fitted model. Given a threshold $q_t$, the prediction set (paper eq.~4) is
\begin{equation}
\C_t(\x)=\{y\in\R:\ s_t(\x,y)\le q_t\}.
\label{eq:cpset}
\end{equation}
The classical choice takes $q_t$ to be the $\lceil(1-\alpha)(t+1)\rceil$-th smallest of the past scores $\{s(\x_\tau,y_\tau)\}_{\tau=1}^t$. Its guarantee rests on exchangeability.

\begin{lemma}[Rank uniformity]\label{lem:rank}
If $S_1,\dots,S_{m+1}$ are exchangeable and a.s.\ distinct, then the rank of $S_{m+1}$ among all $m+1$ values is uniform on $\{1,\dots,m+1\}$; hence for any integer $k$, $\Prob\big(S_{m+1}\le S_{(k)}\big)\ge \tfrac{k}{m+1}$, where $S_{(k)}$ is the $k$-th smallest of $S_1,\dots,S_m$.
\end{lemma}
\begin{proof}
Exchangeability makes all $(m+1)!$ orderings equally likely, so the rank $R$ of $S_{m+1}$ is uniform on $\{1,\dots,m+1\}$. The event $\{S_{m+1}\le S_{(k)}\}$ contains $\{R\le k\}$, and $\Prob(R\le k)=k/(m+1)$.
\end{proof}

\begin{proposition}[Marginal coverage]\label{prop:cov}
With $m=t$ calibration scores and $q_t=S_{(k)}$, $k=\lceil(1-\alpha)(t+1)\rceil$, the set \eqref{eq:cpset} satisfies $\Prob\big(y_{t+1}\in\C_t(\x_{t+1})\big)\ge1-\alpha$.
\end{proposition}
\begin{proof}
$y_{t+1}\in\C_t$ iff $s(\x_{t+1},y_{t+1})=S_{t+1}\le q_t=S_{(k)}$. By Lemma~\ref{lem:rank}, $\Prob(S_{t+1}\le S_{(k)})\ge k/(t+1)\ge1-\alpha$.
\end{proof}

\noindent\textbf{Code: the split-CP threshold.} \code{run\_plain\_cp} returns exactly $S_{(k)}$ with $k=\lceil(1-\alpha)(t+1)\rceil$ (finite-sample correction \code{plus\_one=True}); the \code{min/max} clamp caps $k$ at $r$, returning the largest score rather than the strict $+\infty$ when $(1-\alpha)(r+1)>r$.
\begin{lstlisting}
def run_plain_cp(calib_scores, alpha, plus_one=True):
    calib_scores = np.asarray(calib_scores, float).ravel()
    r = len(calib_scores)
    if r == 0:
        return np.inf
    if plus_one:
        k = int(np.ceil((1 - alpha) * (r + 1)))
    else:
        k = int(np.ceil((1 - alpha) * r))
    k = min(max(k, 1), r)
    return float(np.sort(calib_scores)[k - 1])
\end{lstlisting}

\noindent\textbf{Score used here: the Gaussian NLL.} With the GP predictive \eqref{eq:pred}, the score is the negative predictive log-likelihood
\begin{equation}
s_t(y)=-\log p(y\mid\D_t,\x)=\underbrace{\tfrac12\log(2\pi\sigma_t^2)}_{=:~\text{nll\_min}}+\frac{(y-\hy_t)^2}{2\sigma_t^2}.
\label{eq:nll}
\end{equation}
\code{standard\_cp} is the online form of Proposition~\ref{prop:cov} for this score: at each $t$ it sets $q$ to the empirical $1-\alpha$ quantile of \emph{past} scores, inverts \eqref{eq:nll} to an interval (derived in \S\ref{sec:conf-set}, eq.~\eqref{eq:setinv}), and appends the new score to the calibration memory.
\begin{lstlisting}
def standard_cp(y, mu, var, alpha=0.1):
    T = len(y)
    var = np.maximum(np.asarray(var, float), 1e-12)
    nll_min = 0.5 * np.log(2 * np.pi * var)          # min score, at y = mu
    score = nll_min + (y - mu) ** 2 / (2 * var)      # observed score s_t(Y_t)
    lo = np.full(T, np.nan); hi = np.full(T, np.nan)
    covered = np.zeros(T, bool); size = np.zeros(T); past = []
    for t in range(T):
        q = np.quantile(past, 1-alpha, method="higher") if past else score[t]
        rad = 2 * var[t] * (q - nll_min[t])
        if rad >= 0:
            h = np.sqrt(rad)
            lo[t], hi[t] = mu[t]-h, mu[t]+h
            size[t] = 2 * h
        covered[t] = score[t] <= q
        past.append(score[t])
    return dict(lo=lo, hi=hi, covered=covered, size=size)
\end{lstlisting}
The paper's point: in the stream, calibration and test scores are \emph{not} exchangeable (the model is trained on $\D_t$, so training scores are systematically smaller), and Proposition~\ref{prop:cov} can fail. Sections~\ref{sec:conf-set}--\ref{sec:theorem} restore a long-run guarantee by adapting $q_t$.

%======================================================================
\section{Scalable GP by random Fourier features (\S III.A)}\label{sec:rff}
Assume $\kappa=\sigma_\theta^2\bar\kappa$ with $\bar\kappa$ stationary and normalised, $\bar\kappa(\x,\x)=1$.

\begin{theorem}[Bochner]\label{thm:bochner}
A continuous stationary $\bar\kappa(\x,\x')=\bar\kappa(\x-\x')$ is positive definite iff it is the Fourier transform of a finite nonnegative measure. Normalising that measure to a density $\pi_{\bar\kappa}$ gives (paper eq.~5)
\begin{equation}
\bar\kappa(\x-\x')=\int_{\R^d}\pi_{\bar\kappa}(\mathbf v)\,e^{\,j\mathbf v^\top(\x-\x')}\,d\mathbf v=\E_{\mathbf v\sim\pi_{\bar\kappa}}\!\big[e^{\,j\mathbf v^\top(\x-\x')}\big].
\label{eq:bochner}
\end{equation}
\end{theorem}

Because $\bar\kappa$ is real and even, $\pi_{\bar\kappa}$ is symmetric, so the imaginary part of \eqref{eq:bochner} integrates to zero and
\begin{equation}
\bar\kappa(\x-\x')=\E_{\mathbf v}\!\big[\cos(\mathbf v^\top(\x-\x'))\big]
=\E_{\mathbf v}\!\big[\cos(\mathbf v^\top\x)\cos(\mathbf v^\top\x')+\sin(\mathbf v^\top\x)\sin(\mathbf v^\top\x')\big],
\label{eq:cosexpand}
\end{equation}
using $\cos(a-b)=\cos a\cos b+\sin a\sin b$. Draw $\mathbf v_1,\dots,\mathbf v_D\stackrel{\text{iid}}{\sim}\pi_{\bar\kappa}$ and define the $2D$-dimensional random-feature map (paper eq.~6)
\begin{equation}
\bphi_{\mathbf v}(\x)=\tfrac{1}{\sqrt D}\big[\sin(\mathbf v_1^\top\x),\cos(\mathbf v_1^\top\x),\dots,\sin(\mathbf v_D^\top\x),\cos(\mathbf v_D^\top\x)\big]^\top.
\label{eq:rffmap}
\end{equation}
\begin{proposition}[Unbiasedness]\label{prop:unbiased}
$\E\big[\bphi_{\mathbf v}(\x)^\top\bphi_{\mathbf v}(\x')\big]=\bar\kappa(\x-\x')$, and $\bphi_{\mathbf v}(\x)^\top\bphi_{\mathbf v}(\x')\to\bar\kappa(\x-\x')$ a.s.\ as $D\to\infty$.
\end{proposition}
\begin{proof}
$\bphi_{\mathbf v}(\x)^\top\bphi_{\mathbf v}(\x')=\frac1D\sum_{i=1}^D\big[\sin(\mathbf v_i^\top\x)\sin(\mathbf v_i^\top\x')+\cos(\mathbf v_i^\top\x)\cos(\mathbf v_i^\top\x')\big]$. Each summand has expectation $\bar\kappa(\x-\x')$ by \eqref{eq:cosexpand}; the average is therefore unbiased, and the strong law gives a.s.\ convergence.
\end{proof}

Thus $\bar\kappa(\x,\x')\approx\bphi_{\mathbf v}^\top(\x)\bphi_{\mathbf v}(\x')$, and the parametric surrogate (paper eq.~7) is a Bayesian linear model
\begin{equation}
\check f(\x)=\bphi_{\mathbf v}^\top(\x)\,\btheta,\qquad \btheta\sim\N(\mathbf 0_{2D},\ \sigma_\theta^2\I_{2D}),
\label{eq:param}
\end{equation}
whose induced covariance reproduces the target kernel: $\mathrm{Cov}(\check f(\x),\check f(\x'))=\bphi^\top(\x)\,(\sigma_\theta^2\I)\,\bphi(\x')=\sigma_\theta^2\,\bphi^\top(\x)\bphi(\x')\approx\sigma_\theta^2\bar\kappa=\kappa$.

\noindent\textbf{Code: the feature map \eqref{eq:rffmap}.}
\begin{lstlisting}
def rff_features(X, V):
    """X:(N, d), V:(d, D)  ->  Z:(N, 2D)."""
    D = V.shape[1]
    VX = X @ V
    return np.sqrt(1.0 / D) * np.concatenate([np.sin(VX), np.cos(VX)], axis=1)
\end{lstlisting}

\noindent\textbf{Spectral densities and the samplers.} The frequency law $\pi_{\bar\kappa}$ is the inverse transform of $\bar\kappa$. For the four kernels (1-D scale $\ell$; the code writes \code{length\_scale}$=\ell$):
\begin{center}
\renewcommand{\arraystretch}{1.35}
\begin{tabular}{lll}
\hline
Kernel $\bar\kappa(r)$ & Spectral density $\pi_{\bar\kappa}(\omega)$ & Sampling recipe (code)\\
\hline
RBF $e^{-r^2/2\ell^2}$ & $\N(0,\ell^{-2})$ & $\omega\sim\N(0,1/\ell^2)$\\
Laplace $e^{-|r|/\ell}$ & Cauchy$(0,\ell^{-1})$ & $\omega\sim\mathrm{Cauchy}(0,1/\ell)$\\
Mat\'ern-$\nu$ & Student-$t_{2\nu}$ (scale $\ell^{-1}$) & $\omega=z/\sqrt g,\ z\sim\N(0,\ell^{-2}),\ g\sim\mathrm{Gamma}(\nu,\tfrac1\nu)$\\
Rational quad.\ ($\alpha$) & Gamma mixture of $\N$ & $\omega=z\sqrt u,\ z\sim\N(0,1),\ u\sim\mathrm{Gamma}(\alpha,\tfrac{1}{\alpha\ell^2})$\\
\hline
\end{tabular}
\end{center}
The Mat\'ern and RQ rows use the \emph{Gaussian scale-mixture} representation: if $\omega\mid g\sim\N(0,\tau^2/g)$ and $g\sim\mathrm{Gamma}(\nu,1/\nu)$ (mean $1$), the marginal of $\omega$ is Student-$t$ with $2\nu$ degrees of freedom; mixing $\N(0,u)$ over $u\sim\mathrm{Gamma}$ yields the RQ spectrum. This is why the code draws a per-frequency mixing variable and rescales a Gaussian.
\begin{lstlisting}
def sample_rbf(d, D, length_scale, rng):        # RBF:  omega ~ N(0, 1/l^2)
    return rng.normal(0.0, 1.0 / length_scale, size=(d, D))

def sample_laplacian(d, D, length_scale, rng):  # exp(-|r|/l): omega ~ Cauchy(0, 1/l)
    return cauchy.rvs(loc=0.0, scale=1.0 / length_scale, size=(d, D), random_state=rng)

def sample_matern(d, D, length_scale, nu, rng): # Student-t_{2nu} via Gaussian scale-mixture
    g = gamma.rvs(a=nu, scale=1.0 / nu, size=D, random_state=rng)
    z = rng.normal(0.0, 1.0 / length_scale, size=(d, D))
    return z / np.sqrt(g)[None, :]

def sample_rq(d, D, length_scale, alpha, rng):  # RQ = scale-mixture of RBFs
    u = gamma.rvs(a=alpha, scale=1.0 / (alpha * length_scale ** 2), size=D, random_state=rng)
    z = rng.normal(0.0, 1.0, size=(d, D))
    return z * np.sqrt(u)[None, :]
\end{lstlisting}

%======================================================================
\section{Weight-space posterior and online predictive (\S III.A)}\label{sec:weight}
Stack $\bphi_\tau=\bphi_{\mathbf v}(\x_\tau)$ into $\bm\Phi_t=[\bphi_1,\dots,\bphi_t]^\top$. Model \eqref{eq:param} with likelihood $y_\tau\mid\btheta\sim\N(\bphi_\tau^\top\btheta,\sn^2)$ is a conjugate Gaussian linear model; the posterior is Gaussian, $p(\btheta\mid\D_t)=\N(\btheta;\hat\btheta_t,\bSigma_t)$.

\begin{proposition}[Batch posterior]\label{prop:batch}
\begin{equation}
\bSigma_t=\Big(\sn^{-2}\bm\Phi_t^\top\bm\Phi_t+\sigma_\theta^{-2}\I\Big)^{-1},\qquad
\hat\btheta_t=\sn^{-2}\bSigma_t\bm\Phi_t^\top\y_t.
\label{eq:batchpost}
\end{equation}
\end{proposition}
\begin{proof}
The log-posterior is $-\tfrac{1}{2\sn^2}\|\y_t-\bm\Phi_t\btheta\|^2-\tfrac{1}{2\sigma_\theta^2}\|\btheta\|^2+\text{const}$. Collecting quadratic and linear terms in $\btheta$ gives precision $\sn^{-2}\bm\Phi_t^\top\bm\Phi_t+\sigma_\theta^{-2}\I$ and linear coefficient $\sn^{-2}\bm\Phi_t^\top\y_t$; completing the square yields \eqref{eq:batchpost}.
\end{proof}

The predictive of $y$ at a test $\x$ integrates $\btheta$ out: with $\bphi=\bphi_{\mathbf v}(\x)$,
\begin{equation}
\hy(\x)=\bphi^\top\hat\btheta_t,\qquad \sigma^2(\x)=\bphi^\top\bSigma_t\bphi+\sn^2,
\label{eq:wpred}
\end{equation}
the extra $\sn^2$ being the observation noise. Equation~\eqref{eq:wpred} at $\x_{t+1}$ is the paper's eq.~8, written $\hy_{t+1|t}=\bphi^\top(\x_{t+1})\hat\btheta_t$ and $\sigma^2_{t+1|t}=\bphi^\top(\x_{t+1})\bSigma_t\bphi(\x_{t+1})+\sn^2$.

\noindent\textbf{Code: batch weight-space solve.} \code{gp\_predict} implements \eqref{eq:batchpost}--\eqref{eq:wpred} with $\bSigma_t=\sn^2\mathbf A^{-1}$, $\mathbf A=\bm\Phi^\top\bm\Phi+(\sn^2/\sigma_\theta^2)\I$: then $\hat\btheta_t=\mathbf A^{-1}\bm\Phi^\top\y$ and $\sigma^2=\sn^2\,\bphi^\top\mathbf A^{-1}\bphi+\sn^2$.
\begin{lstlisting}
def gp_predict(Z_train, y_train, Z_test, sigma_n, sigma_f):
    M = Z_train.shape[1]
    A = Z_train.T @ Z_train + (sigma_n ** 2 / sigma_f ** 2) * np.eye(M)
    w = np.linalg.solve(A, Z_train.T @ y_train)          # posterior mean weights
    mu = Z_test @ w
    AinvZt = np.linalg.solve(A, Z_test.T)                # (M, N_test)
    var = sigma_n ** 2 * np.einsum('ij,ji->i', Z_test, AinvZt) + sigma_n ** 2
    return mu, var
\end{lstlisting}
The driver \code{run\_online\_gp} calls this on a growing training set (it therefore stores all past data --- the non-memoryless variant, used for the mini-batch experiments).

%======================================================================
\section{The memoryless recursive update (\S III.A, eqs.~10--11)}\label{sec:recursive}
Absorbing one datum $(\x_{t+1},y_{t+1})$ updates the posterior by Bayes' rule (paper eq.~10),
$p(\btheta\mid\D_{t+1})\propto p(\btheta\mid\D_t)\,\N(y_{t+1};\bphi^\top\btheta,\sn^2)$, again Gaussian. Write $\bphi=\bphi(\x_{t+1})$.

\begin{theorem}[Recursive update, paper eq.~11]\label{thm:rls}
With $\hy_{t+1|t}=\bphi^\top\hat\btheta_t$ and $\sigma^2_{t+1|t}=\bphi^\top\bSigma_t\bphi+\sn^2$,
\begin{align}
\hat\btheta_{t+1}&=\hat\btheta_t+\sigma_{t+1|t}^{-2}\,\bSigma_t\bphi\,(y_{t+1}-\hy_{t+1|t}),\label{eq:thetaup}\\
\bSigma_{t+1}&=\bSigma_t-\sigma_{t+1|t}^{-2}\,\bSigma_t\bphi\bphi^\top\bSigma_t.\label{eq:sigup}
\end{align}
\end{theorem}
\begin{proof}
The updated precision is $\bSigma_{t+1}^{-1}=\bSigma_t^{-1}+\sn^{-2}\bphi\bphi^\top$ (add one likelihood term to Proposition~\ref{prop:batch}). Apply the Sherman--Morrison identity $(\mathbf M+\mathbf u\mathbf u^\top)^{-1}=\mathbf M^{-1}-\frac{\mathbf M^{-1}\mathbf u\mathbf u^\top\mathbf M^{-1}}{1+\mathbf u^\top\mathbf M^{-1}\mathbf u}$ with $\mathbf M=\bSigma_t^{-1}$, $\mathbf u=\sn^{-1}\bphi$:
\[
\bSigma_{t+1}=\bSigma_t-\frac{\sn^{-2}\bSigma_t\bphi\bphi^\top\bSigma_t}{1+\sn^{-2}\bphi^\top\bSigma_t\bphi}
=\bSigma_t-\frac{\bSigma_t\bphi\bphi^\top\bSigma_t}{\sn^2+\bphi^\top\bSigma_t\bphi}
=\bSigma_t-\frac{\bSigma_t\bphi\bphi^\top\bSigma_t}{\sigma^2_{t+1|t}},
\]
which is \eqref{eq:sigup}. For the mean, the conjugate update is $\hat\btheta_{t+1}=\bSigma_{t+1}\big(\bSigma_t^{-1}\hat\btheta_t+\sn^{-2}\bphi\,y_{t+1}\big)=\hat\btheta_t+\sn^{-2}\bSigma_{t+1}\bphi\,(y_{t+1}-\bphi^\top\hat\btheta_t)$, using $\bSigma_{t+1}\bSigma_t^{-1}=\I-\sn^{-2}\bSigma_{t+1}\bphi\bphi^\top$. Finally
\[
\bSigma_{t+1}\bphi=\bSigma_t\bphi-\frac{\bSigma_t\bphi(\bphi^\top\bSigma_t\bphi)}{\sigma^2_{t+1|t}}=\bSigma_t\bphi\,\frac{\sigma^2_{t+1|t}-\bphi^\top\bSigma_t\bphi}{\sigma^2_{t+1|t}}=\bSigma_t\bphi\,\frac{\sn^2}{\sigma^2_{t+1|t}},
\]
so $\sn^{-2}\bSigma_{t+1}\bphi=\sigma_{t+1|t}^{-2}\bSigma_t\bphi$, giving \eqref{eq:thetaup}.
\end{proof}

Each step costs $O((2D)^2)$ (a rank-one update of $\bSigma\in\R^{2D\times2D}$) and stores no past data. \code{run\_online\_gp\_recursive} is \eqref{eq:thetaup}--\eqref{eq:sigup} verbatim: \code{Sigma} starts at $\sigma_\theta^2\I$ (prior), \code{s2} is $\sigma^2_{t+1|t}$, \code{gain}$=\bSigma_t\bphi/\sigma^2_{t+1|t}$, and the symmetrisation line guards against round-off.
\begin{lstlisting}
def run_online_gp_recursive(X_all, y_all, V, sigma_n, sigma_f, n_init=100):
    M = 2 * V.shape[1]
    theta = np.zeros(M)
    Sigma = (sigma_f ** 2) * np.eye(M)          # prior  theta ~ N(0, sigma_f^2 I)
    Z = rff_features(X_all, V); sn2 = sigma_n ** 2
    preds, truths, varis = [], [], []
    for t in range(len(X_all)):
        phi = Z[t]
        mu = phi @ theta
        Sphi = Sigma @ phi
        s2 = float(phi @ Sphi + sn2)            # predictive variance sigma^2_{t+1|t}
        if t >= n_init:
            preds.append(mu); truths.append(y_all[t]); varis.append(s2)
        gain = Sphi / s2                         # recursive update
        theta = theta + gain * (y_all[t] - mu)
        Sigma = Sigma - np.outer(Sphi, Sphi) / s2
        Sigma = 0.5 * (Sigma + Sigma.T)          # keep symmetric
    return np.array(preds), np.array(truths), np.array(varis)
\end{lstlisting}

%======================================================================
\section{Online Bayes credible set (\S III.A, eq.~9)}\label{sec:bayescred}
Using \eqref{eq:wpred} at $\x_{t+1}$, the Bayes $\beta$-credible set is (paper eq.~9)
\begin{equation}
K_t^\beta(\x_{t+1})=[\hy_{t+1|t}-c_\beta\,\sigma_{t+1|t},\ \hy_{t+1|t}+c_\beta\,\sigma_{t+1|t}],
\label{eq:bayescred}
\end{equation}
with $c_\beta$ \emph{constant}. \code{bayes\_credible} uses $c_\beta=\Phi^{-1}(1-\alpha/2)$ (so $\beta=1-\alpha$; $c_\beta\approx1.645$ at $\alpha=0.1$). Under mis-specification the fixed $c_\beta$ makes \eqref{eq:bayescred} mis-covered --- the motivation for a data-driven, time-varying multiplier next.
\begin{lstlisting}
def bayes_credible(y, mu, var, alpha=0.1):
    c = norm.ppf(1 - alpha / 2)                  # fixed multiplier (1.645 @ alpha=0.1)
    sd = np.sqrt(np.maximum(var, 1e-12))
    return dict(lo=mu - c*sd, hi=mu + c*sd,
                covered=np.abs(y - mu) <= c*sd, size=2*c*sd)
\end{lstlisting}

%======================================================================
\section{Conformalised GP set and threshold update (\S III.B)}\label{sec:conf-set}
Take the score \eqref{eq:nll}. Inverting \eqref{eq:cpset} gives a closed interval (paper eq.~12).
\begin{proposition}[Set inversion]\label{prop:setinv}
For $s_t$ in \eqref{eq:nll} and threshold $q_t$,
\begin{equation}
\C_t(\x_{t+1})=\{y:s_t(y)\le q_t\}=\big[\hy_{t+1|t}-c_{t+1|t}\sigma_{t+1|t},\ \hy_{t+1|t}+c_{t+1|t}\sigma_{t+1|t}\big],\quad
c_{t+1|t}=\sqrt{2q_t-\log(2\pi\sigma^2_{t+1|t})},
\label{eq:setinv}
\end{equation}
which is nonempty iff $q_t\ge\tfrac12\log(2\pi\sigma^2_{t+1|t})$, and is always \emph{bounded}.
\end{proposition}
\begin{proof}
$s_t(y)\le q_t\iff \tfrac{(y-\hy)^2}{2\sigma^2}\le q_t-\tfrac12\log(2\pi\sigma^2)\iff(y-\hy)^2\le 2\sigma^2 q_t-\sigma^2\log(2\pi\sigma^2)=\sigma^2(2q_t-\log2\pi\sigma^2)$. The right side is $\ge0$ iff $q_t\ge\tfrac12\log(2\pi\sigma^2)$; taking square roots gives half-width $\sigma\sqrt{2q_t-\log2\pi\sigma^2}=c_{t+1|t}\sigma$. Because $s_t(y)\to\infty$ as $|y|\to\infty$, the sublevel set is bounded (contrast a residual score, whose set can be all of $\R$).
\end{proof}
Unlike the constant $c_\beta$ of \eqref{eq:bayescred}, here $c_{t+1|t}$ "changes over time, adapting to new data samples." The threshold is driven by coverage feedback (paper eq.~13):
\begin{equation}
q_{t+1}=q_t+\eta_t\big(\indic(y_{t+1}\notin\C_t(\x_{t+1}))-\alpha\big).
\label{eq:qupdate}
\end{equation}
\begin{proposition}[eq.~13 is online sub-gradient descent on the pinball loss, eq.~14]\label{prop:pinball}
Let $\rho_{1-\alpha}(u)=(1-\alpha)\max\{u,0\}+\alpha\max\{-u,0\}$ and $u_t=s_t(\x_{t+1},y_{t+1})-q_t$. Then a step $q_{t+1}=q_t-\eta_t\,\partial_{q}\rho_{1-\alpha}(u_t)$ equals \eqref{eq:qupdate}.
\end{proposition}
\begin{proof}
For $u\ne0$, $\rho'_{1-\alpha}(u)=(1-\alpha)$ if $u>0$ and $-\alpha$ if $u<0$. Since $u_t=s_t-q_t$, $\partial_q\rho=-\rho'(u_t)$, i.e.\ $-(1-\alpha)$ when $s_t>q_t$ (a miss, $y_{t+1}\notin\C_t$) and $+\alpha$ when $s_t<q_t$ (covered). Hence $q_{t+1}=q_t-\eta_t\partial_q\rho$ equals $q_t+\eta_t(1-\alpha)$ on a miss and $q_t-\eta_t\alpha$ when covered --- exactly $q_t+\eta_t(\indic(\text{miss})-\alpha)$.
\end{proof}

\noindent\textbf{Code.} \code{online\_conformal} computes the score \eqref{eq:nll}, inverts it (\code{rad}$=2\sigma^2(q-\text{nll\_min})$, \code{h}$=c_{t+1|t}\sigma_{t+1|t}$; empty set when \code{rad}$<0$), records coverage, and applies \eqref{eq:qupdate} in the final line. The \code{step} argument selects $\eta_t$ (\S\ref{sec:schedules}).
\begin{lstlisting}
def online_conformal(y, mu, var, alpha=0.1, step="varying", c=1.0, eps=0.1,
                     W=15, r=100, q1=None):
    y = np.asarray(y, float); T = len(y)
    var = np.maximum(np.asarray(var, float), 1e-12)
    nll_min = 0.5 * np.log(2 * np.pi * var)
    score   = nll_min + (y - mu) ** 2 / (2 * var)
    if q1 is None:
        q1 = float(np.median(nll_min) + chi2.ppf(1.0 - alpha, df=1) / 2)
    q = q1
    lo = np.full(T, np.nan); hi = np.full(T, np.nan)
    size = np.zeros(T); covered = np.zeros(T, bool); q_path = np.empty(T)
    resets, win = [], []; tau = 1; prev_avg = None; rises = 0; a = 0.5 + eps
    for t in range(T):
        q_path[t] = q
        rad = 2 * var[t] * (q - nll_min[t])            # invert the NLL score
        if rad >= 0.0:
            h = np.sqrt(rad)
            lo[t], hi[t] = mu[t] - h, mu[t] + h
            size[t] = 2 * h
        covered[t] = score[t] <= q
        if step == "varying":
            win.append(size[t])
            if len(win) > W: win.pop(0)
            if len(win) == W:                          # set-size change detector
                avg = np.mean(win)
                if (prev_avg is not None and avg > prev_avg): rises += 1
                else: rises = 0
                prev_avg = avg
                if rises >= r:
                    tau = 1; rises = 0; resets.append(t)
            eta = c * tau ** (-a); tau += 1
        elif step == "decaying":
            eta = c * (t + 1) ** (-a)
        else:
            eta = float(step)                          # fixed constant eta
        q = q + eta * ((0.0 if covered[t] else 1.0) - alpha)   # eq. (13)
    return dict(lo=lo, hi=hi, size=size, covered=covered, q_path=q_path, resets=resets)
\end{lstlisting}

%======================================================================
\section{Learning-rate schedules and change-point reset (\S III.B)}\label{sec:schedules}
Three regimes for $\eta_t$, all in the branch above (with $a=\tfrac12+\varepsilon$):
\begin{itemize}[leftmargin=1.4em,itemsep=2pt]
\item \textbf{Constant} $\eta_t=\eta$ (\code{step=<float>}): asymptotic coverage, but the paper notes it "will yield high variability in the instantaneous coverage."
\item \textbf{Decaying} $\eta_t=c\,(t{+}1)^{-a}$ (\code{step="decaying"}): valid under stationary or slow drift; the experiments use $\eta_t=t^{-3/5}$ ($\varepsilon=0.1$).
\item \textbf{Varying} $\eta_t=c\,\tau^{-a}$ with $\tau$ the slots since the last reset (\code{step="varying"}): the set size $|\C_t|=2c_{t+1|t}\sigma_{t+1|t}$ is averaged over a window of $W$ slots, and "if the average prediction set size increases over $r$ consecutive steps, we will declare a distributional shift and then reset the learning rate" ($\tau\leftarrow1$). Defaults $W=15$, $r=100$.
\end{itemize}

%======================================================================
\section{Long-run coverage guarantee (\S III.B, eq.~15)}\label{sec:theorem}
\begin{theorem}[paper eq.~15]\label{thm:cov}
If $s_t:\mathcal X\times\mathcal Y\to[0,B]$ and $q_0\in[0,B]$, then for any data sequence and any positive $\{\eta_t\}$,
\begin{equation}
\left|\frac1T\sum_{t=0}^{T-1}\indic\big(y_{t+1}\in\C_t(\x_{t+1})\big)-(1-\alpha)\right|\le\frac{C\,N_T}{T},\quad
C=\frac{2\big(B+\max_t\eta_t\big)}{\min_t\eta_t},\ \ N_T=\sum_{\tau=1}^{T}\indic(\eta_\tau>\eta_{\tau-1}).
\label{eq:thm}
\end{equation}
\end{theorem}
\begin{proof}[Derivation (constant $\eta$, then the reset correction)]
Let $\mathrm{err}_t=\indic(y_{t+1}\notin\C_t)$, so \eqref{eq:qupdate} is $q_{t+1}=q_t+\eta_t(\mathrm{err}_t-\alpha)$. First take $\eta_t\equiv\eta$. Telescoping,
\[
q_T-q_0=\eta\sum_{t=0}^{T-1}(\mathrm{err}_t-\alpha)\ \Longrightarrow\ \frac1T\sum_{t=0}^{T-1}(\mathrm{err}_t-\alpha)=\frac{q_T-q_0}{\eta T}.
\]
It remains to bound $|q_T-q_0|$. A key monotonicity of the update keeps $q_t$ in a bounded band: if $q_t>B$ then $s_t\le B<q_t$ so the point is covered ($\mathrm{err}_t=0$) and $q_{t+1}=q_t-\eta\alpha\le q_t$; if $q_t<0$ then $s_t\ge0>q_t$ so it is a miss and $q_{t+1}=q_t+\eta(1-\alpha)\ge q_t$. Hence, starting from $q_0\in[0,B]$, $q_t\in[-\eta\alpha,\ B+\eta(1-\alpha)]\subseteq[-\eta,B+\eta]$ for all $t$, so $|q_T-q_0|\le B+\eta$. Therefore
\[
\left|\frac1T\sum_{t=0}^{T-1}(\mathrm{err}_t-\alpha)\right|\le\frac{B+\eta}{\eta T},
\]
and since $\frac1T\sum\indic(y_{t+1}\in\C_t)-(1-\alpha)=-\big(\frac1T\sum\mathrm{err}_t-\alpha\big)$, this is \eqref{eq:thm} with $N_T=1$ (no rate changes). For a general positive schedule, rewrite the telescoped sum with the varying $\eta_t$ and split it at the steps where $\eta_t$ increases: on any block where $\eta_t$ is non-increasing the same clipping argument bounds the block's contribution by $(B+\max_t\eta_t)/\min_t\eta_t$, and there are $N_T$ such increase events, giving the $C\,N_T$ numerator. Full accounting is in Angelopoulos, Barber \& Bates (2024) and Xu et al.\ (2024).
\end{proof}
\noindent The practical reading, quoted from the paper: "as long as the step size does not decay too quickly, and the number of resets $N_T$ is sublinear in $T$, OS-GP-CP can achieve long-run coverage convergence." In code, the returned \code{resets} list \emph{is} the set of increase events counted by $N_T$; the varying schedule aims to keep $|\text{\code{resets}}|=N_T=o(T)$.

%======================================================================
\section{Hyperparameters and experimental setup (\S IV)}\label{sec:hyper}
The kernel magnitude $\sigma_\theta^2$, lengthscale $\ell^2$ and noise $\sn^2$ are fit by maximising the GP log marginal likelihood on the first $100$ points,
\begin{equation}
\log p(\y\mid\mathbf X,\vartheta)=-\tfrac12\y^\top(\K+\sn^2\I)^{-1}\y-\tfrac12\log|\K+\sn^2\I|-\tfrac{t}{2}\log2\pi,
\label{eq:mll}
\end{equation}
then frozen; $D=200$ spectral features and an RBF kernel are used. \code{fit\_hyperparameters} optimises \eqref{eq:mll} through scikit-learn's \code{ConstantKernel*RBF + WhiteKernel} and returns $(\ell,\sigma_\theta,\sn)$.
\begin{lstlisting}
def fit_hyperparameters(X, y, n0=100):
    k = (ConstantKernel(1.0,(1e-3,1e3))*RBF(1.0,(1e-2,1e2)) + WhiteKernel(0.1,(1e-6,1e1)))
    gpr = GaussianProcessRegressor(kernel=k, alpha=0.0, n_restarts_optimizer=2).fit(X[:n0], y[:n0])
    p = gpr.kernel_.get_params()
    return (float(p["k1__k2__length_scale"]),               # length_scale
            float(np.sqrt(p["k1__k1__constant_value"])),    # sigma_f  (= sigma_theta)
            float(np.sqrt(max(p["k2__noise_level"], 1e-6))))# sigma_n
\end{lstlisting}

%======================================================================
\section{Extension in the code: the kernel ensemble}\label{sec:ensemble}
The manuscript's experiments use a single RBF kernel. The code additionally forms a four-kernel ensemble by causal Bayesian model averaging (BMA); this is an add-on, not part of the paper's stated method. Given per-model Gaussian predictives $\N(\mu_i,\sigma_i^2)$ with weights $w_i$, the mixture's first two moments follow from the law of total variance,
\begin{equation}
\mu_{\text{ens}}=\sum_i w_i\mu_i,\qquad
\sigma^2_{\text{ens}}=\sum_i w_i(\sigma_i^2+\mu_i^2)-\Big(\sum_i w_i\mu_i\Big)^2,
\label{eq:mix}
\end{equation}
and the BMA weights update multiplicatively by the marginal likelihood, $w_i\leftarrow w_i\,p_i(y_t)/\!\sum_j w_j p_j(y_t)$, done in log-space (weights for $t$ use only points $<t$, so there is no leakage). \code{gaussian\_loglik} supplies $\log p_i(y_t)$ and \code{ensemble} carries \eqref{eq:mix} plus the weight recursion; the moment-matched $(\mu_{\text{ens}},\sigma^2_{\text{ens}})$ then feeds the conformal layer of \S\ref{sec:conf-set}.
\begin{lstlisting}
def gaussian_loglik(y, mu, var):
    var = np.maximum(var, 1e-12)
    return -0.5 * (np.log(2 * np.pi * var) + (y - mu) ** 2 / var)

def ensemble(preds, varis, logliks, init_weights=None):
    n_models, T = preds.shape
    w = (np.full(n_models, 1.0 / n_models) if init_weights is None
         else np.asarray(init_weights, float))
    ens_mu = np.zeros(T); ens_var = np.zeros(T); weight_hist = np.zeros((T, n_models))
    for t in range(T):
        weight_hist[t] = w
        mu_t = w @ preds[:, t]                              # mixture mean
        second = w @ (varis[:, t] + preds[:, t] ** 2)       # E[y^2] under the mixture
        ens_mu[t] = mu_t; ens_var[t] = second - mu_t ** 2   # law of total variance
        logw = np.log(w + 1e-300) + logliks[:, t]           # BMA weight update (log-space)
        logw -= logw.max(); w = np.exp(logw); w /= w.sum()
    return ens_mu, ens_var, weight_hist
\end{lstlisting}

%======================================================================
\section{Equation-to-code map}\label{sec:map}
\begin{center}
\renewcommand{\arraystretch}{1.3}
\begin{tabular}{lll}
\hline
Paper & Object & Function\\
\hline
(1)--(3) & exact GP joint / predictive & Lemma~\ref{lem:cond}; realised via features (\S\ref{sec:weight})\\
(4) & conformal set + split-CP quantile & \code{run\_plain\_cp}, \code{standard\_cp}\\
(5) & spectral density / Bochner & \code{sample\_rbf/laplacian/matern/rq}\\
(6) & random-feature map $\bphi_{\mathbf v}$ & \code{rff\_features}\\
(7) & parametric model + weight prior & \code{run\_online\_gp\_recursive} (prior $\sigma_\theta^2\I$)\\
(8) & online predictive mean/variance & \code{run\_online\_gp\_recursive}, \code{gp\_predict}\\
(9) & Bayes credible set & \code{bayes\_credible}\\
(10)--(11) & recursive weight-posterior update & \code{run\_online\_gp\_recursive} (Thm.~\ref{thm:rls})\\
(12) & conformalised set, $c_{t+1|t}$ & \code{online\_conformal} (Prop.~\ref{prop:setinv})\\
(13)--(14) & threshold feedback update & \code{online\_conformal} (Prop.~\ref{prop:pinball})\\
schedules & constant/decaying/varying $\eta_t$ & \code{online\_conformal} (\code{step}, \code{resets})\\
(15) & coverage theorem, $N_T$ & Thm.~\ref{thm:cov}; \code{resets} are the $N_T$ events\\
\S IV & marginal-likelihood fit & \code{fit\_hyperparameters}\\
--- & kernel ensemble (extension) & \code{gaussian\_loglik}, \code{ensemble}\\
\hline
\end{tabular}
\end{center}

\end{document}
